\section{Noncommutativity}

\begin{exercise}[An entangled state has no local state vector]
\textbf{Experiment.}
Two-qubit tomography reconstructs joint states whose correlations cannot be
represented by assigning a pure state to each qubit.  Let
\(\rho\) be the reconstructed positive trace-one operator on
\(H_A\otimes H_B\).

Derive the unique operator \(\rho_A\) satisfying
\[
        \Tr(\rho_A X)=\Tr\!\bigl(\rho(X\otimes\id)\bigr)
\]
for every observable \(X\) on \(H_A\).  For a maximally entangled pure state,
show that \(\rho_A\) is maximally mixed.  Explain how a globally sharp state
can therefore give locally random results.  Show that no choice of bases in
\(H_A\) and \(H_B\) turns this state into a product, and identify the
basis-free invariant that detects this failure.
\end{exercise}

\begin{exercise}[An ancilla improves discrimination]
\textbf{Experiment.}
Entangling a probe with an untouched ancilla can increase the probability of
distinguishing two laboratory channels.  For channels \(\Phi_0,\Phi_1\), the
largest unassisted bias is governed by their action on system states, while
the assisted bias allows states on \(H\otimes A\).

Express the assisted bias as the norm of
\((\Phi_0-\Phi_1)\otimes\id_A\).  Show that an ancilla of dimension
\(\dim H\) is sufficient.  Construct a pair of channels for which the
assisted value is strictly larger than the unassisted value.  Explain why the
experiment forces a matrix-level norm on maps rather than a norm which tests
only the original system.
\end{exercise}

\begin{exercise}[Bell correlations]
\textbf{Experiment.}
Loophole-reduced Bell tests observe correlations violating the CHSH bound
\(2\) while remaining below the quantum bound \(2\sqrt2\)
\cite{hensen}.  Let \(A_0,A_1\) and \(B_0,B_1\) be self-adjoint contractions
belonging to commuting local observable algebras, and set
\[
        C=A_0(B_0+B_1)+A_1(B_0-B_1).
\]

Prove the classical bound when all four observables commute.  In the quantum
case, compute \(C^2\) and derive \(\|C\|\le2\sqrt2\)
\cite{tsirelson}.  Give observables and a bipartite state attaining equality.
Explain which algebraic change permits the experimentally observed violation
and why the same calculation still forbids correlations beyond the quantum
bound.
\end{exercise}

\begin{exercise}[Two photons leave together]
\textbf{Experiment.}
When two indistinguishable photons enter opposite ports of a balanced beam
splitter simultaneously, coincident detections at the two outputs are
suppressed \cite{hong-ou-mandel}.  Let \(H\) be the one-photon mode space and
describe two identical bosons in the symmetric square of \(H\).

Apply the beam-splitter unitary functorially to the two-photon state.  Show
that the two amplitudes leading to one photon in each output cancel, whereas
the amplitudes for both photons in one output add.  Explain why distinguishable
photons do not show complete suppression.  Relate the width and depth of the
observed coincidence dip to temporal overlap and indistinguishability rather
than to a choice of mode coordinates.
\end{exercise}

\begin{exercise}[Sequential spin measurements]
\textbf{Experiment.}
A beam selected by a vertical Stern--Gerlach analyzer becomes uncertain when
measured horizontally; selecting the horizontal output then restores two
possible vertical outcomes.  Represent the selections by rank-one
projections \(P\) and \(Q\).

Compute the sequential probabilities from products such as \(PQP\).  Show
that the observed order effect disappears exactly when the projections
commute.  Explain why the disturbance is not accounted for by ignorance of a
pre-existing pair of spin components.  Express the uncertainty through the
commutator without selecting matrices until the final numerical example.
\end{exercise}

\begin{exercise}[No perfect photocopier]
\textbf{Experiment.}
Optical experiments can approximately clone states from a known ensemble, but
no device produces two perfect copies of every unknown polarization state.
Suppose a physical evolution \(U\) and a fixed blank state \(b\) satisfied
\[
        U(v\otimes b)=v\otimes v
\]
for every unit vector \(v\).

Compare inner products before and after the evolution and derive a
contradiction for two distinct nonorthogonal states.  Explain why known
orthogonal states can be copied and why probabilistic or approximate cloning
does not contradict the argument.  Identify preservation of transition
probabilities as the structure forced by the experiment.
\end{exercise}

\begin{exercise}[Teleportation correlations]
\textbf{Experiment.}
Quantum teleportation transfers an unknown qubit state using a shared
entangled pair, two classical bits, and no transmission of the original
carrier.  Let \(H\) be the qubit space and use the canonical correspondence
between maps on \(H\) and vectors in \(H\otimes H^*\).

Decompose the combined state into four orthogonal measurement alternatives.
Show that each alternative leaves the remote system related to the input by a
known unitary correction.  Explain why averaging before receipt of the
classical message gives a maximally mixed remote state, preventing
faster-than-light signaling.  State what experimental statistics certify
fidelity beyond a classical measure-and-prepare scheme.
\end{exercise}

\begin{exercise}[Decoherence of interference]
\textbf{Experiment.}
Interference visibility falls when path information leaks into an
environment, even if no observer reads that information.  Let the two paths
become correlated with normalized environment states \(e_0,e_1\).

Take the partial trace of the joint pure state and show that the
off-diagonal interference term is multiplied by
\(\langle e_1,e_0\rangle\).  Relate fringe visibility to this overlap and path
distinguishability to its complement.  Explain the disappearance of fringes
when the records are orthogonal and their recovery under a quantum eraser
which recombines the record alternatives.
\end{exercise}

\begin{exercise}[Quantum Zeno inhibition]
\textbf{Experiment.}
Repeated measurements of whether an unstable or coherently driven system
remains in its initial state can inhibit its transition.  Let \(P\) project
onto that state and let \(U(t)\) be the unobserved evolution.

Show that after \(N\) equally spaced ideal tests the survival probability is
\[
        \bigl\|(PU(t/N)P)^Nv\bigr\|^2.
\]
Use the short-time quadratic behavior to find its limit as \(N\to\infty\).
Explain why the result depends on the noncommuting alternation of evolution
and projection, and state which finite measurement times limit the laboratory
effect.
\end{exercise}
